\documentclass[a4paper,12pt]{article}

\usepackage[T2A]{fontenc}
\usepackage[utf8]{inputenc}
\usepackage[russian]{babel}
\usepackage{amsmath,amssymb,mathtools}
\usepackage{PTSans}
\renewcommand{\familydefault}{\sfdefault}
\usepackage{icomma}
\usepackage[
    left=20mm,
    right=20mm,
    top=20mm,
    bottom=22mm
]{geometry}
\usepackage{microtype}
\usepackage{tikz}
\usetikzlibrary{calc,arrows.meta,positioning,shapes.geometric}
\usepackage{tkz-euclide}
\usepackage{pgfplots}
\pgfplotsset{compat=1.18}
\usepackage{tabularx,booktabs}
\usepackage{enumitem}
\usepackage{parskip}
\usepackage{fancyhdr}
\usepackage{setspace}
\usepackage{xcolor}

\definecolor{accent}{HTML}{2A6F73}
\definecolor{darkaccent}{HTML}{17484B}
\definecolor{textgray}{HTML}{333333}
\definecolor{softgray}{HTML}{E8EEEE}
\definecolor{warning}{HTML}{8A4B35}

\color{textgray}
\onehalfspacing
\setlength{\parindent}{0pt}
\setlength{\headheight}{25pt}

\setlist[itemize]{
    leftmargin=6mm,
    itemsep=2pt,
    topsep=3pt
}

\setlist[enumerate]{
    leftmargin=8mm,
    itemsep=4pt,
    topsep=4pt
}

\pagestyle{fancy}
\fancyhf{}
\fancyhead[L]{\small\color{darkaccent}Чек-лист по математике}
\fancyhead[R]{\small\color{darkaccent}27.09.26}
\fancyfoot[C]{\small\thepage}
\renewcommand{\headrulewidth}{0.4pt}
\renewcommand{\headrule}{
    \hbox to\headwidth{
        \color{softgray}
        \leaders\hrule height \headrulewidth\hfill
    }
}

\newcommand{\keyidea}[1]{%
    \medskip
    {\color{darkaccent}\bfseries #1}
    \smallskip
}

\newcommand{\important}[1]{%
    {\color{warning}\bfseries #1}
}

\begin{document}

% ============================================================
% ТИТУЛЬНЫЙ ЛИСТ
% ============================================================

\begin{titlepage}
\thispagestyle{empty}

\vspace*{24mm}

\begin{center}

{\Large\color{darkaccent}\bfseries ЧЕК-ЛИСТ}

\vspace{9mm}

{\huge\bfseries
Сложные задачи ОГЭ:\\[3mm]
неравенства, графики, движение\\
и подобие треугольников}

\vspace{12mm}

{\large
Повторение ключевых методов занятия
}

\vspace{22mm}

\begin{tabular}{rl}
\textbf{Ученица:} & София Кальней \\[3mm]
\textbf{Преподаватель:} & Артём Александрович Лёвин \\[3mm]
\textbf{Дата:} & 27.09.26
\end{tabular}

\vfill

{\large\color{darkaccent}
Лёвин Артём Александрович\\
эксклюзивно для Софии Кальней
}

\vspace{18mm}

\begin{tikzpicture}[x=1cm,y=1cm]
\draw[
    accent,
    line width=1.4pt
]
(0,0)
.. controls (3,0.8) and (6,-0.8) .. (9,0)
.. controls (11.5,0.8) and (14,-0.4) .. (16.4,0.25);
\end{tikzpicture}

\end{center}

\end{titlepage}

% ============================================================
% 1. НЕРАВЕНСТВА
% ============================================================

\section*{\color{darkaccent}
1. Сложные неравенства: сначала ищите структуру
}

Перед механическим раскрытием скобок проверьте, можно ли преобразовать
выражение короче.

\keyidea{Разность квадратов}

Если в неравенстве встречается сравнение квадратов
\[
A^2\ge B^2,
\]
удобно перенести всё в одну сторону:
\[
A^2-B^2\ge0,
\]
а затем использовать
\[
A^2-B^2=(A-B)(A+B).
\]

Это часто позволяет избежать громоздкого раскрытия квадратов.

\subsection*{Пример}

Решим:
\[
(9x-4)^2\ge(4x-9)^2.
\]

Переносим правую часть:
\[
(9x-4)^2-(4x-9)^2\ge0.
\]

Разность квадратов:
\[
\bigl((9x-4)-(4x-9)\bigr)
\bigl((9x-4)+(4x-9)\bigr)\ge0.
\]

Получаем
\[
(5x+5)(13x-13)\ge0,
\]
то есть
\[
65(x+1)(x-1)\ge0.
\]

Поскольку \(65>0\),
\[
(x+1)(x-1)\ge0.
\]

Критические точки:
\[
x=-1,\qquad x=1.
\]

Ответ:
\[
\boxed{x\in(-\infty;-1]\cup[1;+\infty)}.
\]

\keyidea{Почему этот способ особенно важен при больших степенях}

Рассмотрим:
\[
(2x^2-1)^2\le(x^2-2)^2.
\]

Если раскрывать оба квадрата напрямую, появится четвёртая степень.
Разность квадратов позволяет этого избежать:
\[
(2x^2-1)^2-(x^2-2)^2\le0.
\]

Получаем:
\[
(x^2+1)(3x^2-3)\le0,
\]
или
\[
3(x^2+1)(x^2-1)\le0.
\]

Но
\[
x^2+1>0
\]
при любом действительном \(x\). Поэтому этот множитель на знак
произведения не влияет:
\[
x^2-1\le0.
\]

Отсюда
\[
\boxed{x\in[-1;1]}.
\]

\subsection*{Ключевой навык: оценивать знак множителя}

Если Вы видите выражение
\[
x^4+x^{100}+47,
\]
то
\[
x^4\ge0,\qquad
x^{100}\ge0,\qquad
47>0,
\]
поэтому
\[
\boxed{x^4+x^{100}+47>0}
\]
при любом действительном \(x\).

Такой множитель можно безопасно исключить делением:
знак неравенства сохраняется.

\important{
Нельзя делить неравенство на выражение с переменной,
пока его знак и возможность обращения в нуль не установлены.
}

Перед делением проверьте:

\begin{itemize}
    \item выражение не равно нулю на рассматриваемых значениях;
    \item его знак известен;
    \item при делении на отрицательное выражение знак неравенства меняется.
\end{itemize}

\subsection*{Чётные степени}

Для любого действительного \(a\):
\[
a^{2n}\ge0.
\]

Кроме того,
\[
(-a)^{2n}=a^{2n}.
\]

Поэтому, например,
\[
(3x-1)^4=(1-3x)^4.
\]

% ============================================================
% 2. НЕОТРИЦАТЕЛЬНЫЕ ВЫРАЖЕНИЯ
% ============================================================

\section*{\color{darkaccent}
2. Сумма неотрицательных выражений равна нулю
}

Если
\[
A\ge0,\qquad B\ge0,
\]
то
\[
A+B=0
\]
возможно только тогда, когда
\[
\boxed{A=0\quad\text{и}\quad B=0}.
\]

Этот приём работает для квадратов, модулей, чётных степеней
и арифметических корней чётной степени.

\subsection*{Пример с модулями}

Решим:
\[
|x^2-x-6|+|x^2+4x+4|=0.
\]

Оба слагаемых неотрицательны, поэтому необходимо одновременно:
\[
\begin{cases}
x^2-x-6=0,\\
x^2+4x+4=0.
\end{cases}
\]

Разложим:
\[
x^2-x-6=(x-3)(x+2),
\]
\[
x^2+4x+4=(x+2)^2.
\]

Общий корень двух уравнений:
\[
\boxed{x=-2}.
\]

\keyidea{Что необходимо замечать}

При виде суммы квадратов, модулей или других гарантированно
неотрицательных выражений сначала выясните, когда каждое слагаемое
может обратиться в нуль.

% ============================================================
% БЛОК-СХЕМА
% ============================================================

\clearpage
\thispagestyle{plain}

\begin{center}

\begin{tikzpicture}[
    node distance=8mm,
    every node/.style={
        font=\small,
        align=center
    },
    title/.style={
        font=\Large\bfseries,
        text=darkaccent,
        align=center
    },
    start/.style={
        ellipse,
        draw=accent,
        line width=1pt,
        minimum width=44mm,
        minimum height=10mm
    },
    action/.style={
        rectangle,
        rounded corners=2mm,
        draw=accent,
        line width=1pt,
        text width=50mm,
        minimum height=11mm
    },
    decision/.style={
        diamond,
        draw=darkaccent,
        line width=1pt,
        aspect=2.7,
        text width=34mm,
        inner sep=1.3mm
    },
    arr/.style={
        -{Stealth[length=3mm]},
        line width=1pt,
        color=darkaccent
    }
]

\node[title] (title)
{Алгоритм: как упростить\\сложное неравенство};

\node[start,below=8mm of title] (s)
{Получено сложное неравенство};

\node[action,below=8mm of s] (move)
{Перенесите всё в одну сторону и сравните выражение с нулём};

\node[decision,below=11mm of move] (square)
{Есть ли структура \(A^2-B^2\)?};

\node[action,below left=12mm and 3mm of square] (factor)
{Разложите:\\
\(A^2-B^2=(A-B)(A+B)\)};

\node[action,below right=12mm and 3mm of square] (simplify)
{Ищите другое полезное преобразование или разложение};

\node[decision,below=34mm of square] (sign)
{Есть множитель с заранее известным знаком?};

\node[action,below left=12mm and 3mm of sign] (remove)
{Если он строго положителен, его можно исключить делением};

\node[action,below right=12mm and 3mm of sign] (keep)
{Если знак неизвестен, делить на него нельзя};

\node[action,below=34mm of sign] (interval)
{Найдите нули оставшихся множителей и примените метод интервалов};

\node[start,below=8mm of interval] (finish)
{Запишите ответ с учётом строгих и нестрогих границ};

\draw[arr] (s)--(move);
\draw[arr] (move)--(square);

\draw[arr]
(square)--node[above left]{да}(factor);

\draw[arr]
(square)--node[above right]{нет}(simplify);

\draw[arr] (factor)|-(sign);
\draw[arr] (simplify)|-(sign);

\draw[arr]
(sign)--node[above left]{да}(remove);

\draw[arr]
(sign)--node[above right]{нет}(keep);

\draw[arr] (remove)|-(interval);
\draw[arr] (keep)|-(interval);

\draw[arr] (interval)--(finish);

\end{tikzpicture}

\end{center}

\clearpage

% ============================================================
% 3. СРЕДНЯЯ СКОРОСТЬ
% ============================================================

\section*{\color{darkaccent}3. Средняя скорость}

Формула средней скорости:
\[
\boxed{
v_{\text{ср}}
=
\frac{S_{\text{общ}}}{t_{\text{общ}}}
}.
\]

\subsection*{Равные половины пути}

Пусть первую половину пути тело прошло со скоростью \(v_1\),
вторую --- со скоростью \(v_2\).

Обозначим каждую половину пути через \(S\). Тогда
\[
S_{\text{общ}}=2S,
\qquad
t_{\text{общ}}
=
\frac{S}{v_1}+\frac{S}{v_2}.
\]

Следовательно,
\[
v_{\text{ср}}
=
\frac{2S}
{\frac{S}{v_1}+\frac{S}{v_2}}
=
\boxed{
\frac{2v_1v_2}{v_1+v_2}
}.
\]

\subsection*{Пример}

Для скоростей \(66\) км/ч и \(78\) км/ч:
\[
v_{\text{ср}}
=
\frac{2\cdot66\cdot78}{66+78}
=
\frac{10296}{144}
=
\boxed{71{,}5\text{ км/ч}}.
\]

% ============================================================
% 4. МНОГОЭТАЖНЫЕ ДРОБИ
% ============================================================

\section*{\color{darkaccent}4. Многоэтажные дроби}

Полезные формы:
\[
\boxed{
\frac{\frac ab}{\frac cd}
=
\frac{ad}{bc}
},
\qquad
\boxed{
\frac{\frac ab}{c}
=
\frac{a}{bc}
},
\qquad
\boxed{
\frac{a}{\frac bc}
=
\frac{ac}{b}
}.
\]

В четырёхэтажной дроби удобно помнить:
внешние множители переходят в числитель,
внутренние --- в знаменатель.

Перед сокращением проверяйте, какие выражения могут обращаться в нуль.

% ============================================================
% 5. ГРАФИК
% ============================================================

\section*{\color{darkaccent}
5. Рациональная функция и выколотая точка
}

Рассмотрим:
\[
y=
\frac{(x-3)(x^2+6x+8)}{x+2}.
\]

\important{Первый шаг --- ОДЗ:}
\[
x\ne-2.
\]

Разложим квадратный трёхчлен:
\[
x^2+6x+8=(x+2)(x+4).
\]

При \(x\ne-2\):
\[
y=
\frac{(x-3)(x+2)(x+4)}{x+2}
=
(x-3)(x+4).
\]

Получаем
\[
y=x^2+x-12,
\qquad
x\ne-2.
\]

График совпадает с параболой
\[
y=x^2+x-12,
\]
но точка при \(x=-2\) удалена.

Её ордината по упрощённой формуле:
\[
y(-2)=4-2-12=-10.
\]

Выколотая точка:
\[
\boxed{(-2;-10)}.
\]

\subsection*{Вершина параболы}

\[
x_{\text{в}}
=
-\frac{b}{2a}
=
-\frac12,
\]
\[
y_{\text{в}}
=
\left(-\frac12\right)^2
-\frac12
-12
=
-\frac{49}{4}.
\]

Вершина:
\[
\boxed{
\left(
-\frac12;
-\frac{49}{4}
\right)
}.
\]

\begin{center}

\begin{tikzpicture}

\begin{axis}[
    width=0.78\linewidth,
    height=7.2cm,
    axis lines=middle,
    xmin=-5,
    xmax=4,
    ymin=-14,
    ymax=8,
    samples=250,
    grid=both,
    minor tick num=1,
    xlabel={$x$},
    ylabel={$y$},
    clip=false
]

\addplot[
    accent,
    line width=1.4pt,
    domain=-5:4
]
{x^2+x-12};

\addplot[
    only marks,
    mark=o,
    mark size=4pt,
    line width=1.2pt,
    darkaccent
]
coordinates {(-2,-10)};

\addplot[
    only marks,
    mark=*,
    mark size=2.5pt,
    darkaccent
]
coordinates {(-0.5,-12.25)};

\node[
    anchor=south west,
    font=\small,
    darkaccent
]
at (axis cs:-2,-10)
{$(-2;-10)$};

\node[
    anchor=north west,
    font=\small,
    darkaccent
]
at (axis cs:-0.5,-12.25)
{$\left(-\frac12;-\frac{49}{4}\right)$};

\end{axis}

\end{tikzpicture}

\end{center}

\subsection*{
Когда прямая \(y=m\) имеет ровно одну общую точку с графиком?
}

Есть два случая.

\begin{enumerate}

    \item Прямая проходит через вершину параболы:
    \[
    m=-\frac{49}{4}.
    \]

    \item Прямая
    \[
    y=-10
    \]
    пересекала бы полную параболу в двух точках,
    но одна из них --- удалённая точка
    \[
    (-2;-10).
    \]
    Поэтому остаётся только одно пересечение.

\end{enumerate}

Итак,
\[
\boxed{
m=-\frac{49}{4}
\quad\text{или}\quad
m=-10
}.
\]

\keyidea{Алгоритм для сокращаемой рациональной функции}

\begin{enumerate}
    \item Выпишите ОДЗ исходной функции.
    \item Разложите числитель и знаменатель на множители.
    \item Сократите общие множители с учётом ОДЗ.
    \item Постройте график упрощённой функции.
    \item Вернитесь к ограничениям исходной функции.
    \item Удалите запрещённые точки и подпишите их координаты.
\end{enumerate}

% ============================================================
% 6. ПОДОБИЕ
% ============================================================

\section*{\color{darkaccent}
6. Геометрия: как увидеть подобие
}

Если в задаче есть несколько треугольников,
попробуйте связать их через подобие.

Сначала ищите равные углы:

\begin{itemize}
    \item общий угол;
    \item вертикальные углы;
    \item углы, равные вследствие параллельности прямых.
\end{itemize}

После этого проверяйте признаки подобия.

\keyidea{Признак по двум сторонам и углу между ними}

Пусть
\[
AB=10,\qquad BC=14,
\]
а точки \(F\in AB\) и \(E\in BC\) удовлетворяют
\[
BF=5,\qquad BE=7.
\]

Тогда
\[
\frac{BF}{BA}
=
\frac{5}{10}
=
\frac12,
\qquad
\frac{BE}{BC}
=
\frac{7}{14}
=
\frac12.
\]

Угол при вершине \(B\) общий:
\[
\angle FBE=\angle ABC.
\]

Следовательно,
\[
\triangle BFE\sim\triangle BAC.
\]

Коэффициент подобия малого треугольника к большому:
\[
k=\frac12.
\]

Если
\[
AC=15,
\]
то
\[
EF=
\frac12\cdot15
=
\boxed{7{,}5}.
\]

\begin{center}

\begin{tikzpicture}[scale=0.48]

\tkzDefPoint(0,0){B}
\tkzDefPoint(10,0){A}
\tkzDefPoint(3.55,13.54){C}

\tkzDefMidPoint(B,A)
\tkzGetPoint{F}

\tkzDefMidPoint(B,C)
\tkzGetPoint{E}

\tkzDrawPolygon[
    line width=1.2pt,
    color=accent
](A,B,C)

\tkzDrawSegment[
    line width=1pt,
    color=darkaccent
](F,E)

\tkzDrawPoints[
    color=darkaccent,
    fill=darkaccent,
    size=3
](A,B,C,E,F)

\tkzLabelPoints[below](B,F,A)
\tkzLabelPoints[left](E)
\tkzLabelPoints[above](C)

\end{tikzpicture}

\end{center}

\keyidea{Трапеция и пересечение диагоналей}

В трапеции \(ABCD\), где
\[
AD\parallel BC,
\]
диагонали \(AC\) и \(BD\) пересекаются в точке \(O\).

Треугольники
\[
\triangle AOD
\quad\text{и}\quad
\triangle COB
\]
подобны по двум углам, поэтому
\[
\boxed{
\frac{AO}{OC}
=
\frac{AD}{BC}
}.
\]

Если
\[
AD=16,\qquad
BC=9,\qquad
AC=15,
\]
обозначим
\[
AO=x,
\qquad
OC=15-x.
\]

Из подобия:
\[
\frac{x}{15-x}
=
\frac{16}{9}.
\]

Отсюда
\[
9x=240-16x,
\]
\[
25x=240,
\]
\[
\boxed{AO=9{,}6}.
\]

% ============================================================
% 7. ИТОГ
% ============================================================

\section*{\color{darkaccent}
7. Что важно уметь после занятия
}

\begin{itemize}

    \item[$\square$]
    Я замечаю разность квадратов до раскрытия скобок.

    \item[$\square$]
    Я умею оценивать знак целого множителя.

    \item[$\square$]
    Я не делю неравенство на выражение неизвестного знака.

    \item[$\square$]
    Я понимаю, когда сумма неотрицательных выражений равна нулю.

    \item[$\square$]
    Я использую полное расстояние и полное время
    при поиске средней скорости.

    \item[$\square$]
    Я умею работать с многоэтажными дробями.

    \item[$\square$]
    Я сначала выписываю ОДЗ рациональной функции.

    \item[$\square$]
    После сокращения функции я возвращаю выколотую точку.

    \item[$\square$]
    В геометрии я ищу общие, вертикальные
    и связанные с параллельностью углы.

    \item[$\square$]
    После доказательства подобия я правильно сопоставляю
    сходственные стороны.

\end{itemize}

% ============================================================
% ТРЕНИРОВОЧНЫЙ БЛОК
% ============================================================

\clearpage

\section*{\color{darkaccent}
Тренировочный блок --- Путь А
}

Решите 10 задач.
Решения оформляйте самостоятельно.

\begin{enumerate}

\item
Решите неравенство:
\[
(5x-2)^2\ge(2x-5)^2.
\]

\item
Решите неравенство:
\[
(3x^2-2)^2\le(x^2-4)^2.
\]

\item
Решите неравенство:
\[
(x^8+x^4+9)(2x-5)\le0.
\]

\item
Решите уравнение:
\[
|x^2-5x+6|+|x^2-4x+4|=0.
\]

\item
Решите неравенство:
\[
(2x-1)^4\ge(2x+1)^4.
\]

\item
Первую половину пути автомобиль прошёл
со скоростью \(60\) км/ч,
вторую половину --- со скоростью \(90\) км/ч.

Найдите среднюю скорость на всём пути.

\item
Вычислите:
\[
\frac{\frac56}{\frac{10}{9}}.
\]

\item
Дана функция
\[
y=
\frac{(x-1)(x^2+5x+6)}{x+2}.
\]

Найдите все значения \(m\), при которых прямая
\[
y=m
\]
имеет с графиком этой функции ровно одну общую точку.

\item
В треугольнике \(ABC\)
\[
AB=15,\qquad
BC=10,\qquad
AC=20.
\]

Точка \(F\) лежит на \(AB\),
точка \(E\) --- на \(BC\), причём
\[
BF=6,\qquad
BE=4.
\]

Найдите \(EF\).

\item
В трапеции \(ABCD\)
\[
AD\parallel BC.
\]

Диагонали пересекаются в точке \(O\).
Известно:
\[
AD=15,\qquad
BC=10,\qquad
AC=25.
\]

Найдите \(AO\).

\end{enumerate}

% ============================================================
% ОТВЕТЫ
% ============================================================

\clearpage

\section*{\color{darkaccent}Ответы}

\renewcommand{\arraystretch}{1.5}

\begin{center}

\begin{tabularx}{\linewidth}{@{}cX@{}}

\toprule
\textbf{№} & \textbf{Ответ} \\
\midrule

1 &
\(
x\in(-\infty;-1]\cup[1;+\infty)
\)
\\

2 &
\(
x\in
\left[
-\sqrt{\frac32};
\sqrt{\frac32}
\right]
\)
\\

3 &
\(
x\in
\left(
-\infty;
\frac52
\right]
\)
\\

4 &
\(
x=2
\)
\\

5 &
\(
x\le0
\)
\\

6 &
\(
72\text{ км/ч}
\)
\\

7 &
\(
\frac34
\)
\\

8 &
\(
m=-4,\,-3
\)
\\

9 &
\(
EF=8
\)
\\

10 &
\(
AO=15
\)
\\

\bottomrule

\end{tabularx}

\end{center}

\vfill

\begin{center}
{\small\color{darkaccent}
Лёвин Артём Александрович эксклюзивно для Софии Кальней
}
\end{center}

\end{document}